% ======================================================================
\chapter{Conclusão}
\label{cap:conclusao}
% ======================================================================
\instrucao{Uma a duas páginas. Retome cada objetivo específico e diga se foi alcançado, com os números principais. Feche com limitações e sugestões concretas para trabalhos futuros.}

O tubo de ressonância permitiu medir a velocidade do som no ar com incerteza relativa abaixo de \qty{0.5}{\percent}, usando apenas uma trena e um gerador de sinais. O ajuste linear das seis primeiras ressonâncias separou o comprimento de onda da correção de extremidade, que de outro modo seria uma fonte de erro sistemático.

Como continuação, sugere-se repetir a medida em várias frequências, para verificar que $v$ não depende de $f$, e em várias temperaturas, para testar a Equação~\eqref{eq:temperatura}.
